\changecolours{ScoutGreen}
\section[Referências]{Referências Bibliográficas Recomendadas}

% ------------------------------------------------------------------------------
\begin{frame}{Referências Bibliográficas}
  \itemseps[topsepi=0.5mm,itemsepi=1mm,labelsep=2mm,leftmargini=3.5mm]
  \begin{itemize}
    \scriptsize
    \item \textbf{HEUSER, Carlos Alberto.} \textit{Projeto de Banco de Dados}. 6. ed. Porto Alegre: Bookman, 2009. (Capítulos 2 e 3: Notação Conceitual no padrão brModelo e Dicionário de Dados).
    \item \textbf{CHEN, Peter Pin-Shan.} \textit{The Entity-Relationship Model: Toward a Unified View of Data}. ACM Transactions on Database Systems (TODS), New York, v. 1, n. 1, p. 9--36, Mar. 1976.
    \item \textbf{ELMASRI, Ramez; NAVATHE, Shamkant B.} \textit{Sistemas de Banco de Dados}. 7. ed. São Paulo: Pearson Education do Brasil, 2018. (Capítulos 3 e 4: Modelagem Conceitual e EER).
    \item \textbf{SILBERSCHATZ, Abraham; KORTH, Henry F.; SUDARSHAN, S.} \textit{Sistema de Banco de Dados}. 7. ed. Rio de Janeiro: Elsevier, 2020.
    \item \textbf{CANDIDO, Carlos H.} \textit{brModelo Web \& Desktop: Ferramenta CASE para Modelagem de Bancos de Dados}. Disponível em: \texttt{https://www.brmodeloweb.com}.
  \end{itemize}
  \vspace{0.15cm}
  \begin{center}
    \tiny\color{gray}
    Instituto Federal de Educação, Ciência e Tecnologia do Ceará -- Campus Tauá\\
    Curso Superior de Tecnologia em Telemática
  \end{center}
\end{frame}
